\subsection*{一个日志位置先表达一项候选}

先把业务去重暂时放在一边，看看副本怎样决定要应用哪条命令。设 N1、N2、N3 的共同日志都已到位置 40，业务状态也已应用至 40。账户 1 有 10000 分，购买、订单、流水和逐笔绑定为空。登记会话 E1 有效。N1 是当前领导者，任期为 7，收到 S1 的 3000 分购买命令。

Raft 把领导者选举的一轮逻辑时期称为任期，用递增整数标识。任期 7 表示这批协议消息属于哪一轮领导权，消息携带它，接收方据此识别落后消息。第八章 T1 的 \code{lease\_epoch} 比较同一恢复任务的领取资格，Raft 任期约束一个复制组中的选举与日志交换，两者分别标记各自协议的代次~\cite{raft}。

N1 把命令写为日志项 \code{(index=41,term=7,command=S1及原输入)}。\code{index} 是日志位置，决定这项候选在共同命令序列中的先后。\code{term} 是该项由领导者接收并创建时的任期。相同 S1 稍后可以在另一个位置出现，应用层到时再按完整身份判断是否已经处理。

N1 把新日志项发送给 N2、N3。某副本保存了 41，说明它已经保存一项候选。领导者按照提交规则确定该项进入持续保留的前缀，才允许状态机应用它。状态机实际执行完，才能得到此次购买的 P/O/D 和原结果。结果发回客户端，又是一个通信步骤~\cite{raft}。

\CElevenFigure{ch11-original04}{接收日志、复制、提交、应用与客户端获知依次形成不同事实}{fig:ch11-five-facts}

设 N1 和 N2 已持久保存 41，N3 尚未收到。在固定三投票副本的配置里，N1 拿到当前任期日志项的两份保存依据，可以提交它。随后 N1 应用 41，余额成为 7000，并保存 S1 的结果。N3 可以稍后追上。这条路径允许少数副本慢一些，同时保留已经达到提交条件的前缀。

再把故障放在 N1 单独写完 41、N2 尚未接收的位置。N1 失效，N2、N3 可以选出新领导者，它们未必保留这条候选。客户端超时以后重发 S1，新领导者可以在自己的合法日志位置重新受理。应用从提交前缀执行，避免把尚可能被替换的候选提前做成业务事实。

\subsection*{相同前驱怎样约束日志}

让日志保持一致，需要检查追加位置以前的历史。N1 发送新项时同时携带 \code{prevLogIndex} 和 \code{prevLogTerm}，指出新项以前那个位置及其任期。N2 只有在自己的该位置具有相同任期时才接受这次追加。缺少前驱或任期不同，就拒绝，让领导者向前寻找共同位置，再发送后缀~\cite{raft}。

另起一条独立历史，三个副本共同已应用的前缀到 39。N1、N2 在位置 40 保存同一条任期 6 的日志项，N3 的位置 40 则是任期 5 的未提交冲突项。N1 当选任期 7 的领导者以后，在位置 41 追加 S1。N2 的 40 与 N1 相同，可以接收 41。N3 的 40 与请求携带的前驱任期不同，拒绝直接接入 41。N1 先找到两者共同的 39，再用自己的 40、41 替换 N3 的未提交冲突后缀。

\CElevenFigure{ch11-original05}{追加请求检查前驱的位置与任期，从共同位置修复冲突后缀}{fig:ch11-log-matching}

这个检查支持一个归纳关系。初始空日志显然共享相同前缀。一位合法领导者在一个任期内为某个位置创建唯一日志项，并且只向自己的日志尾部追加。跟随者接受新项以前先匹配前驱，所以每次成功追加都把相同前缀向后扩展一段。由此得到日志匹配性质：两个副本在相同位置保存相同任期的项时，该项内容以及以前的全部前缀一致。

跟随者可能拥有额外的未提交项，修复过程可以删除这些冲突项。领导者自己在当选后的任期内保持追加，选举规则则要保证它从一开始就带着过去已经提交的前缀。否则，一位缺少旧成功的领导者会把其他副本的正确历史当成冲突后缀替掉。

\subsection*{为什么投票还要比较日志}

Raft 的投票者在同一任期最多投一票，任期和投票决定在响应以前持久保存。三副本中的任意两个多数集合至少有一个交集成员。因此同一任期的两个候选者无法各取得两票：那个交集成员需要投两票才能让两者都成功。进程重启时保留自己投过谁，才能延续这项约束~\cite{raft}。

同任期领导者唯一以后，要保住旧提交，投票者还比较候选者的日志是否至少与自己一样新。比较顺序先看最后一项的任期，再在任期相同时比较最后位置。末项为 \code{(term=8,index=42)} 的日志比 \code{(term=7,index=100)} 更新。这里按日志产生的任期判断，任期相同时才用长度打破平局。

\CElevenFigure{ch11-vote-order}{选票先比较末项任期，任期相同时再比较末项位置}{fig:ch11-vote-order}

先看一个直观情况。N1 在任期 7 把 41 提交到 N1/N2，N1 随后失效。N3 只有到 40 的日志，向 N2 请求任期 8 的票。N2 的末项任期 7、位置 41，比 N3 更新，于是拒绝。N2 带着 41 参加选举，可以得到 N3 的票，成为新领导者。新的多数与旧的提交多数通过 N2 接上同一项事实。

一般证明还要覆盖候选者拥有更高末项任期的情况。设任期 T 的领导者提交了自己的日志项 e，假定后来出现第一位缺少 e 的领导者，其任期为 U。旧提交多数与 U 的投票多数有一位共同成员 V。V 必须先接收 e，再投票给 U，因为它一旦进入 U，就会拒绝来自较小任期 T 的追加。

在 U 以前的领导者都保存 e，这是选择第一位缺少 e 的领导者所给的条件。领导者保持追加，跟随者只接受这些领导者的匹配前缀，所以 V 投票时继续保存 e。接下来分两种情况。若 U 候选者的末项任期与 V 相同，它获得票就要求日志至少一样长。相同末项的任期及其前缀匹配关系使它包含 e。若它的末项任期更大，那条末项来自 U 以前的一位领导者。该领导者包含 e，候选者持有它的末项，也就带着包含 e 的前缀。两种情况都与缺少 e 矛盾~\cite{raft}。

\CElevenFigure{ch11-original06}{提交多数与后继选举多数的交集，把已经提交的前缀传到新领导者}{fig:ch11-leader-completeness}

由此得到领导者完整性：一个任期中提交的当前任期日志项，会出现在所有更高任期领导者的日志中。状态机只按已提交项的日志顺序应用，各副本在同一个位置就会应用同一命令。不同副本可以在不同时间追赶到这个位置，已经应用的命令内容保持一致。

\subsection*{旧任期项为什么不能只数副本}

前面的提交条件特意说当前任期。任期变化以后，单纯看到旧项落在多数副本上，有时不足以判定它安全。Raft 原文用五个副本构造了这个反例~\cite{raft}。下面把每个位置的内容补全，以便逐步检查。

设五副本 A、B、C、D、E 的位置 1 都相同。任期 2 的领导者 A 在位置 2 写 x，复制到 A、B 后失效，x 尚未达到三份。任期 3，E 的日志只有共同位置 1。它得到 C、D、E 的票，在自己的位置 2 写 y，y 暂时只留在 E，随后 E 失效。

任期 4，A 恢复，带着任期 2 的 x，获得 A、B、C 的票成为领导者。它把 x 追加到 C。现在 A、B、C 都有位置 2 的 x，确实共有三份。如果此时就把这个旧项判为提交，接下来的合法选举会暴露矛盾。

A 又失效，E 恢复，参加任期 5 的选举。E 的末项任期为 3，B、C 的末项任期为 2，D 只有共同位置 1。因此 E 可以得到自己以及 B、C 的票。它成为领导者以后，在 B、C 上用任期 3 的 y 替换位置 2 的 x。x 曾经在三台机器上出现，却没有满足保证后继领导者继承它的提交条件。

\CElevenFigure{ch11-original07}{五副本反例：旧任期 x 达到多数以后，更高末项任期的 y 可以赢得选举}{fig:ch11-old-term}

修复发生在任期 4。A 在复制 x 后，再于位置 3 写入一个任期 4 的项 z，并将它复制到 A、B、C。z 可以是一条不修改业务数据的空命令。当前任期项 z 达到多数，按照前面证明的完整性关系提交。它的前驱 x 随着同一前缀共同提交。E 末项任期 3，已经无法从具有任期 4 末项的 B、C 取得足够票。

因此领导者推进 \code{commitIndex} 时，选一个得到多数保存依据、且项任期等于当前任期的位置 N。\code{commitIndex} 表示该副本目前知道已经提交的最高日志位置。把它推进到 N，也就覆盖以前的全部项。协议在旧项达到更多副本以后，可以继续复制并通过当前任期项确定整个前缀。应用程序依赖库给出的提交边界执行业务转换。

\subsection*{已经提交，怎样变成可恢复的结果}

\code{lastApplied} 表示本副本已经应用到的最高位置。它可以暂时小于 \code{commitIndex}。例如领导者已经知道到 45 的前缀提交，本地状态机只执行到 43，S1 恰好位于 44。此时直接查询应用映射会得到 S1 不存在，但日志中已有等待应用的受理事实。先让状态机追到相应边界，再解释查询，才能把两个层次接起来。

\CElevenFigure{ch11-commit-apply}{提交边界在 45，应用边界在 43，S1 的绑定将在应用 44 后出现}{fig:ch11-commit-apply}

安全响应的一条清楚路径是把请求提交为命令，等待该项提交并且在本地状态机应用，然后返回应用所得结果。领导者在这之后、发回结果以前失效，客户端可能超时。下一节会让重发命令从已经应用的绑定取得原结果。

只读查询也有相同的新鲜度问题。一个与多数隔离的旧领导者可能尚未听说新任期，自己的状态也落后。仅检查本地领导者标志（如 \code{role=leader}），就直接返回 S1 缺失或旧余额，会违反第六章的线性化读取要求~\cite{raft,ongaro}。

一种容易理解的办法，是让只读查询也进入有序日志，等到它提交并按序执行。Raft 还提供避免为每次读取写日志的办法：新领导者先确认自己任期中的一项已经提交，通常由空命令完成。读取开始后记录当前提交边界 r，通过新一轮与多数副本的联系确认本轮领导资格，再等待本地状态机至少应用到 r，最后读取对应应用状态。\code{readIndex} 就是这里用于等待的读取边界~\cite{ongaro}。

\CElevenFigure{ch11-original08}{线性化读取同时确认本轮领导资格与应用位置，去重判断另行核对 S 和原输入}{fig:ch11-read-boundary}

多数联系应与这一批读取关联。跟随者提供读取时，也可以先取得领导者确认的读取边界，再等待自己应用到那里。用租约减少这轮联系，则要把时钟误差与资格期限纳入安全前提。第八章讨论的租约条件在这个具体读协议上再次出现。

线性化读确定这次读取可以放在哪个有效时间点，S 绑定确定两条命令是否属于同一次购买。一张正确的 S 去重表也可能被落后读取路径读成空。随后若命令重新进入权威状态机，应用仲裁可以恢复原结果。若把读空直接当成跨系统执行授权，旧付款就可能被重复发起。读取协议先建立有效观察，操作仲裁再按完整身份解释观察所得状态。

\subsection*{安全和进展分别依赖什么}

上述安全推导采用固定投票成员、可信协议参与者和崩溃恢复模型。网络可以延迟、丢失、重复或重排消息，分区期间少数一侧无法取得提交多数。节点在响应以前保存任期、投票和日志，恢复后沿已保存状态继续。磁盘无声损坏、恢复早期备份或人为建立另一套同名新集群，会改变这些保存前提。节点恶意伪造投票和内容属于另一类故障模型~\cite{raft,paxos-made-simple}。

进展要求存在足够长的运行区间：多数副本与其存储可用，消息交换和持久化通常能在选举超时以前完成，领导者有机会持续复制，客户端能够找到有效入口并重试。Raft 用随机化超时降低持续分票的机会，时延持续大于超时则会反复选举~\cite{raft}。第二章介绍的 FLP 结论已经说明，完全异步模型里，即使只容许一个进程停止，确定性共识也无法为所有执行保证终止~\cite{flp}。

\CElevenFigure{ch11-safety-progress}{前缀保留、业务提交、应用追赶与客户端获知分别依赖对应条件}{fig:ch11-safety-progress}

把这些条件放回当前例子，多数前缀能够延续，说明 S1 所在命令可以继续被恢复。应用追赶以后，响应通道与持续重试使客户端有机会取得原结果。外部平台的 X1 则按第七章的查询与执行条件确定，再由本地恢复协议保存证据。
\FloatBarrier
